\documentclass{ximera}

\input{../../preamble.tex}

\title{Chaining Functions}

\begin{document}

\begin{abstract}
composition
\end{abstract}
\maketitle





Using two given functions, we can create a third function, called the \textbf{composition}.

If $f$ and $g$ are the names of the two component functions, then $f \circ g$ is the symbol for the composition of $f$ and $g$.

We compose $f$ and $g$ to form $(f \circ g)$.

The order makes a difference.

We compose $g$ and $f$ to form $(g \circ f)$.


Most of time these are not equal functions.

\[
(f \circ g) \ne (g \circ f)
\]


\begin{idea} \textbf{\textcolor{blue!55!black}{Composition}}


Let $F$ and $G$ be two functions.

The \textbf{composition of $F$ and $G$} is another function called the composition of $F$ and $G$.  A circle between the function names denotes a composition.

\[
F \circ G
\]


\[
(F \circ G)(a) = F(G(a))
\]


The composition applies $G$ to an element of the domain of $G$ and then uses the resulting function value from $G$ as a domain member of $F$ obtaining a function value of $F$.

This ultimate value of $F$ resulting from this process is the value of $F \circ G$.

\end{idea}



Composition is our way of representing an ordered application of processes.



\subsection*{Algebraic Composition}


Our first version of algebraic representations for functions is formulas (there will be other versions).

Each component function has its own formula, using its own variable, which may or may not represent the numbers in the domain.

The composition is evaluated by first evaluating the inner function at the domain number and then evaluating the outer function at the value of the inner function.




\textbf{Two Component Functions:}


Let $H$ and $K$ be two functions defined by the following formulas.

\[
H(n) = 2 n^2 - 5 \,  \text{ with domain }  \,  (-3, 8)
\]


\[
K(v) = -v + 1 \,  \text{ with domain }  \,  (-10, 5)
\]





\begin{example}

From the formula for $H$, we can see that $H(0) = -5$.

From the formula for $K$, we can see that $K(-5) = 6$.

These two values tell us that $(K \circ H)(0) = K(H(0)) = K(-5) = 6$

\end{example}






\begin{example}

From the formula for $K$, we can see that $K(2) = -1$.

From the formula for $H$, we can see that $H(-1) = -3$.

These two values tell us that $(H \circ K)(2) = H(K(2)) = H(-1) = -3$

\end{example}









\begin{example}

From the formula for $H$, we can see that $H(-3) = 13$.

From the formula for $K$, we can see that $K(13) = -12$.

These two values tell us that $(K \circ H)(-3) = K(H(-3)) = K(13) = -12$

\end{example}






\begin{example}

From the formula for $K$, we can see that $K(-9) = 10$.

From the formula for $H$, we can see that $H(10)$ is undefined.

We cannot include $10$ in the domain of $(H \circ K)$.

\end{example}










\begin{question}




Let $A$ and $B$ be two functions defined by the following formulas.

\[
A(x) = 3 - 2x \,  \text{ with domain }  \,  (-5, 5]
\]


\[
B(y) = -y^2 \,  \text{ with domain }  \,  [-9, 5]
\]




From the formula for $B$, we can see that $B(1) = \answer{-1}$.

From the formula for $A$, we can see that $A(-1) = \answer{1}$.

These two values tell us that $(A \circ B)(1) = \answer{1}$

\end{question}








\begin{question}




Let $M$ and $T$ be two functions defined by the following formulas.

\[
M(x) = 3 - 2x \,  \text{ with domain }  \,  (-5, 5]
\]


\[
T(y) = -y^2 \,  \text{ with domain }  \,  [-9, 5]
\]





From the formula for $M$, we can see that $M(3) = \answer{-3}$.

From the formula for $T$, we can see that $T(-3) = \answer{-9}$.

These two values tell us that $(T \circ M)(3) = \answer{-9}$

\end{question}









\begin{question}




Let $R$ and $N$ be two functions defined by the following formulas.

\[
R(t) = 3 - 2t \,  \text{ with domain }  \,  (-5, 5]
\]


\[
N(k) = -k^2 \,  \text{ with domain }  \,  [-9, 5]
\]


Is $4$ in the domain of $(N \circ R)$?

\begin{multipleChoice}
\choice [correct]{Yes}
\choice {No}
\end{multipleChoice}







Is $4$ in the domain of $(R \circ N)$?

\begin{multipleChoice}
\choice {Yes}
\choice [correct]{No}
\end{multipleChoice}



\end{question}








\begin{onlineOnly}
\begin{center}
\textbf{\textcolor{green!50!black}{ooooo-=-=-=-ooOoo-=-=-=-ooooo}} \\

more examples can be found by following this link\\ \link[More Examples of Formulas]{https://ximera.osu.edu/csccmathematics/precalculus1/precalculus1/formulas/examples/exampleList}

\end{center}
\end{onlineOnly}





\end{document}
